\section{Learning the Allocation with Abstention}
\label{sec:learning}

The oracle search requires per-video codec round-trip and SAM2 inference---infeasible at dataset scale.
We ask: \emph{can a lightweight predictor learn to approximate the oracle allocation from first-frame features alone?}

\subsection{Experimental Setup}

We train the two-stage predictor (\cref{subsec:predictor}) on oracle-labelled DAVIS videos.
Training uses MSE loss on projected sector alphas plus BCE gate loss ($\lambda{=}0.5$), 300 epochs of Adam ($\mathrm{lr}{=}10^{-3}$).
Evaluation is real H.264 + SAM2 inference on \emph{held-out} videos (never seen during training).

\paragraph{Cross-validation protocol.}
We use 5-fold random CV on $n{=}28$ oracle-labelled videos (22 train / 6 test per split), with fresh model training per split.
This is a conservative evaluation: the predictor must generalise to unseen video content, object categories, and motion patterns.

\subsection{Results}

\begin{table}[t]
\centering
\small
\caption{Learned predictor vs.\ uniform BRS (5-fold CV, $n{=}28$, 30 held-out evaluations total). ``Closure'' $= (\text{learned} - \text{BRS}) / (\text{oracle} - \text{BRS})$.}
\label{tab:predictor}
\begin{tabular}{lccc}
\toprule
\textbf{Metric} & \textbf{One-stage} & \textbf{Two-stage (gate)} & \textbf{$\Delta$} \\
\midrule
Win-rate vs.\ uniform    & 83\%  & \textbf{90\%}  & $+7$ \\
Mean gain (pp)           & 14.9  & \textbf{17.6}  & $+2.7$ \\
Mean closure ($|$gap$|{>}5$pp) & 26.4\% & \textbf{31.4\%} & $+5.0$ \\
Negative losses          & 5/30  & \textbf{3/30}  & $-2$ \\
Min split win-rate       & 66.7\% & \textbf{83.3\%} & $+16.6$ \\
\bottomrule
\end{tabular}
\end{table}

The two-stage predictor achieves $90\%$ win-rate over uniform BRS with $+17.6\,\mathrm{pp}$ mean gain.
The gate reduces catastrophic losses from 5 to 3 out of 30 evaluations, and lifts the worst-split win-rate from $66.7\%$ to $83.3\%$.
Mean closure (the fraction of the oracle gap recovered) is $31.4\%$ for videos with meaningful gap ($|$gap$|{>}5\,\mathrm{pp}$).

\paragraph{Interpretation.}
The predictor learns a real signal---not random search noise---as evidenced by consistent ${\geq}83\%$ win-rate across all 5 random splits.
However, $31\%$ closure with $n{=}22$ training videos and 60-dim hand-crafted features is a proof of concept, not a ceiling.
The gate mechanism is crucial: without it, the predictor occasionally produces allocations worse than uniform BRS on out-of-distribution videos (e.g., unusual object shapes or very small masks).
With the gate, the predictor substantially reduces---but does not eliminate---the risk of underperforming uniform BRS.

\subsection{Failure Modes}

Three videos consistently resist learned allocation across splits:
\textbf{bear} (closure $-13\%$), \textbf{elephant} ($-64\%$), and \textbf{kite-surf} ($-75\%$).
These share a pattern: either the oracle gap is small (uniform BRS already strong, little room for improvement) or the optimal sector pattern requires directional features not captured by first-frame statistics alone (e.g., future motion trajectory).
The gate catches most such cases but not all---future work should incorporate temporal features or SAM2 encoder representations.
